\chapter{Introduction}
\label{chap:introduction}

\section{Background of Study}
\label{sec:1-1-background-of-study}

The growth of social media platforms such as WeChat, Facebook, and Instagram has made social media a major marketing and communication channel \parencite{appel2020,mangold2009}. Beyond their original function of connecting individuals, these platforms have become important marketing channels. Businesses disseminate promotional messages, interact with target customers, and encourage users to participate in the circulation of content. Unlike traditional offline marketing channels, social media allows marketing messages to be shared, liked, commented on, and redistributed through interpersonal connections, allowing a single piece of content to travel across multiple layers of a network and reach audiences beyond the firm's direct contact base.

This networked communication environment has changed both the mechanism and the measurement of marketing effectiveness. In traditional marketing, firms usually deliver messages through controlled channels such as print media, television, outdoor advertisements, or direct sales. In social media marketing, however, marketing effectiveness relates to whether users open a message, spend time reading it, share it with others, and allow it to move through their own social relationships. Such peer-to-peer communication pattern can extend a firm's reach beyond its direct audience, and additionally, digital content marketing can support engagement and trust through relevant brand-related content and consumer interactions \parencite{hollebeek2019,mangold2009}.

One increasingly important form of social media marketing relies on a firm's own internal or agent-based social network. Under this model, a company distributes promotional articles to employees, sales representatives, or marketing agents, who then share the articles through their personal WeChat networks. In this dissertation, these sharers are referred to as marketing agents. They are important because they form the entry points through which firm-generated content enters interpersonal networks. Each agent's initial share exposes an article to that agent's direct contacts, and any subsequent resharing may carry the message further into the surrounding network. The agent is both an organisational representative and the starting point of a diffusion process.

Despite the advantages of agent-based social media marketing, a persistent practical challenge remains: the same marketing message does not diffuse equally well through all agents. An article that attracts attention and resharing when posted by one agent may receive limited engagement when posted by another. This unevenness suggests that diffusion effectiveness relates not only to the general quality of the article itself, but also to the characteristics of the agent who shares it and to whether the article's topic is relevant to the agent's professional context and audience. For example, an article about home decoration or renovation design may be more suitable for an agent whose role, expertise, or client network is closely connected to interior design, while the same article may be less effective when shared through a less relevant agent context.

This issue is closely related to the limitation of one-size-fits-all content distribution. In many social marketing systems, identical articles may be distributed broadly to many agents without sufficient consideration of topic relevance, agent role, or audience context. Such a strategy may appear efficient because it maximises the number of potential sharers, but it can also waste limited social capacity and weaken immediate diffusion performance.

For these reasons, this dissertation focuses on the relationship between article content, observable agent characteristics, and content–agent fit in social media marketing diffusion, where diffusion effectiveness is understood as a set of related outcomes. The study does not treat agents as interchangeable channels for broadcasting the same content. It examines whether marketing messages are associated with stronger diffusion when their topics fit the agent context through which they are shared. The empirical setting is the WeiMai social marketing platform, a marketing service used by client firms to manage agent-based WeChat marketing. The data analysed here are drawn from a single interior-design company that uses this platform, and record real-world WeChat article-sharing and reading behaviour. By analysing how articles, agents, and their alignment relate to observed diffusion outcomes, this dissertation seeks to provide evidence for moving beyond uniform content broadcasting towards more personalised and context-sensitive content assignment.

\section{Research Problem}
\label{sec:1-2-research-problem}

Agent-based social media marketing provides firms with an opportunity to disseminate promotional content through interpersonal networks. The effectiveness of this approach, however, remains uncertain because marketing messages do not diffuse uniformly across agents, audiences, or content categories. In practice, a firm may hold a large collection of marketing articles and a pool of agents who can share these articles through their personal WeChat networks. The managerial challenge is not simply to produce attractive content, but to determine which articles should be assigned to which agents in ways that support stronger diffusion effectiveness.

Existing studies on online diffusion have generally emphasised two broad sets of factors: message content and social network structure. Content-oriented research suggests that the characteristics of a message, such as its topic, relevance, emotional appeal, and presentation format, may be associated with users' willingness to read and reshare it. Network-oriented research, by contrast, suggests that the sender's position and the surrounding network context may relate to the scale and shape of diffusion. Both perspectives are important, but each remains incomplete when considered alone. A content-only explanation cannot fully account for why the same article performs differently when shared by different agents, while an agent-only explanation cannot fully account for why the same agent generates different diffusion outcomes for different types of content.

The central problem addressed in this dissertation is that marketing diffusion effectiveness may be associated not only with article content or agent characteristics separately, but also with whether a particular article topic fits the agent context through which it is shared. The research problem can be summarised as follows: how are article content, observable agent characteristics, and content–agent fit associated with the multi-dimensional diffusion effectiveness of marketing messages in an agent-based WeChat marketing system?

In this study, the agent is treated as both the origin point of diffusion and an observable representation of a surrounding interpersonal network. Agent characteristics, such as professional role and demographic attributes, are used as practical proxies for the local audience context and potential content preferences embedded in the agent's network. This correspondence between what is communicated and who communicates it is referred to in this dissertation as content–agent fit.

A further challenge is how to define diffusion effectiveness itself. In social marketing, effectiveness cannot be fully captured by a single reach measure. A message may reach many readers but remain shallow, or it may reach fewer readers while generating deeper cascades and more resharing. This dissertation therefore treats diffusion effectiveness as a multi-dimensional concept. Reconstructed cascade and network metrics are used as outcome variables or supplementary descriptors only where their level-specific meaning is coherent.

This issue also creates an important level-of-analysis problem. Content characteristics, agent characteristics, and content–agent fit are not observed at the same empirical level. Content effects are most naturally examined across article–agent diffusion cases, agent characteristics are examined across agents, and content–agent fit is examined across agent–topic combinations. Combining these heterogeneous units into one undifferentiated model would make coefficient interpretation unclear and could increase the risk of treating outcome-derived diffusion metrics as if they were independent explanatory variables. The analysis adopts a three-level design that separates content-level article–agent cases, agent-characteristic-level records, and content–agent matching-level agent–topic pairs.

Addressing this problem has both academic and practical significance. Academically, the study develops a leakage-aware empirical framework that separates content, agent characteristics, and matching while placing reconstructed diffusion-network measures on the outcome side when their construction requires it. Practically, the study provides evidence for moving beyond one-size-fits-all content distribution. If stronger content–agent fit is associated with stronger diffusion performance, firms may improve social marketing efficiency by assigning articles more selectively to agents whose professional contexts are better aligned with the content topic.

\section{Research Questions}
\label{sec:1-3-research-questions}

Building on the research problem identified above, this dissertation examines whether article content, agent characteristics, and content–agent fit are systematically associated with the diffusion effectiveness of marketing messages. The study is not designed to establish causal effects, nor is it framed as a predictive comparison among alternative machine-learning models. Instead, it adopts an explanatory, associational approach organised around three units of analysis, so that each research question corresponds to a clear unit. Diffusion effectiveness is understood as a multi-dimensional outcome. Reach-type measures serve as the primary outcomes at each level, while cascade-shape and network-position measures are used as supplementary outcomes.

The dissertation is guided by the following three research questions.

\noindent \textbf{RQ1: Are article content characteristics associated with article–agent diffusion effectiveness in WeChat marketing?}

The first research question concerns the content level. It examines whether observable article characteristics are associated with diffusion outcomes when articles are shared through agents' WeChat networks. These characteristics include topic category, title–content semantic consistency, word count, and image-use indicators that are all derived from the article files.

\noindent \textbf{RQ2: Are observable agent characteristics associated with agent-level diffusion outcomes?}

The second research question concerns the agent-characteristic level. Because each agent is both an embedded member of the social network and the origin point from which diffusion begins, an agent's observable attributes can be read as proxies for the local audience and sharing context that the agent represents. This question examines whether agents differ in their average diffusion outcomes and whether such differences are associated with observable agent characteristics such as job category and gender, while accounting for the agent's activity level. Consistent with the multi-dimensional view of diffusion effectiveness, the reconstructed cascade and network metrics computed at this level, such as cascade depth, structural virality, the Wiener index, and centrality, are examined as agent-level outcomes or descriptors.

\noindent \textbf{RQ3: Is content–agent fit associated with stronger agent–topic diffusion performance?}

The third research question concerns the content–agent matching level and represents the central concern of the dissertation. It examines whether stronger alignment between an article's topics and the sharing agent's professional context is associated with stronger diffusion performance at the agent–topic level. This question directly addresses the one-size-fits-all problem by asking whether the same type of content may perform differently depending on the agent context through which it is shared, and it provides the most direct test of the dissertation's central proposition that, in agent-based social marketing, one size does not fit all.

These questions correspond to the three empirical levels used in the analysis: article–agent cases, agents, and agent–topic pairs.

\section{Research Contributions}
\label{sec:1-4-research-contributions}

First, the study provides empirical evidence from a real agent-based WeChat marketing system. The analysis uses behavioural records generated through the WeiMai social marketing platform, giving more process detail than simulated diffusion data or platform-level aggregate statistics. These records capture how promotional articles are shared by marketing agents and subsequently read or reshared through WeChat interpersonal networks. This empirical setting makes it possible to examine marketing diffusion as it occurs in a practical firm-side context, where organisational members act as the entry points through which content enters personal social networks. By focusing on a single interior-design company, the dissertation provides context-specific evidence on social marketing diffusion in an industry where professional relevance, trust, and topic fit are likely to shape marketing communication.

Second, the dissertation contributes an LLM-assisted approach for transforming unstructured marketing articles into structured content features. The article files contain titles, body text, and image-use information, which are difficult to incorporate directly into conventional empirical models. To address this issue, the study uses large language model–assisted content analysis to construct topic labels, topic probability scores, image indicators, word count, and title–content semantic consistency measures. This process allows unstructured article content to be represented in a quantitative form suitable for statistical analysis while preserving the substantive meaning of article topics and message coherence.

Third, the dissertation develops a leakage-aware three-level empirical framework for analysing social marketing diffusion. Instead of combining all variables into one oversized master model, the study separates the analysis into three distinct units: article–agent diffusion cases, agent-level records, and agent–topic pairs. This structure allows the content model, the agent-characteristic model, and the content–agent matching model to answer different research questions with clearer coefficient interpretation. It also reduces the risk of mixing variables that operate at different observational levels.

Fourth, the study clarifies the role of reconstructed cascade and diffusion-network measures in empirical diffusion analysis. Measures such as diffusion depth, structural virality, the Wiener index, and centrality describe useful aspects of observed diffusion, but because they are reconstructed from realised cascades they should not be used as pre-diffusion predictors of the same process. Placing these metrics on the outcome side contributes to a more careful empirical treatment of diffusion data and helps avoid outcome-derived information leakage.

Fifth, the dissertation provides practical implications for personalised content assignment in social media marketing. The study is motivated by the limitation of one-size-fits-all content distribution, where identical articles are distributed broadly to agents without sufficient consideration of topic relevance, agent role, or audience context. By examining whether content–agent fit is associated with stronger diffusion performance, the dissertation offers evidence that may help firms assign articles more selectively to agents whose professional contexts are better aligned with the content topic. Such an approach can support more efficient use of agents' limited social capacity and provide decision support for social marketing recommendation systems.
